\documentclass[12pt,a4paper]{article}
\usepackage[utf8]{inputenc}
\usepackage{amsmath, graphicx, hyperref, listings, booktabs}
\usepackage[margin=2.5cm]{geometry}
\title{Assignment 4: Final Capstone Project — Reservoir Dispatch Optimization}
\author{Hikmatullah Danish \quad 3125999089}
\date{July 2026}

\begin{document}
\maketitle

\section{Introduction}
This capstone project integrates all skills from the AI-Augmented Software Engineering course into a complete mini-application for reservoir dispatch optimization. The system uses Sequential Least Squares Programming (SLSQP) to solve a seven-day drought-period dispatch that balances hydropower revenue against ecological flow requirements. It demonstrates the full development pipeline: prompt engineering, test-driven development, physical validation, multi-objective Pareto analysis, and systematic documentation.

\section{Problem Formulation}
The reservoir has an initial storage of 500,000~m$^3$, bounded by 100,000--1,000,000~m$^3$. Release rates must stay between the ecological minimum of 10~m$^3$/s and the maximum of 100~m$^3$/s. The objective is to maximize revenue:
\begin{equation}
\max \quad \text{Revenue} = K_{\text{hydro}} \sum_{t=1}^{7} Q_t \cdot 24 \cdot p_t
\end{equation}
subject to mass balance $V_{t+1} = V_t + (I_t - Q_t)\cdot 86400$. Six constraints are verified in the validation suite.

\section{Architecture}
The application follows a package structure with \texttt{reservoir\_capstone.py} containing all modules: core simulation (\texttt{compute\_storage}), revenue calculation (\texttt{compute\_total\_revenue}), ecological deficit (\texttt{compute\_eco\_deficit}), SLSQP optimizer, multi-objective Pareto generator, visualization, and validation/export. Outputs are organized into \texttt{outputs/}: \texttt{optimal\_schedule.csv}, \texttt{pareto\_frontier.png}, and \texttt{validation\_report.txt}.

\section{Multi-Objective Analysis}
A Pareto frontier is generated by sweeping 20 weight pairs $(w_{\text{rev}}, w_{\text{eco}})$ in a weighted-sum scalarization. Both objectives are normalized by their maximum possible values to ensure comparable weight influence. The frontier reveals the tangible cost of ecological protection.

\section{Testing \& Validation}
The pytest suite contains 10 tests covering mass balance, revenue positivity, revenue monotonicity, ecological deficit conditions, and output file integrity. A separate validation script runs 13 automated checks on the optimal solution: release bounds, storage bounds, mass balance, revenue consistency, ecological deficit, file existence, and numerical stability. All tests and checks pass.

\section{Results}
The optimal schedule achieves \$54,916.25 in total revenue with zero ecological deficit. The reservoir respects all 6 physical constraints throughout the 7-day horizon. The Pareto frontier confirms that revenue decreases as ecological weight increases, quantifying the cost of environmental protection.

\section{Conclusion}
This capstone demonstrates that a complete optimization application—from problem formulation through validated solution—can be developed using iterative AI-assisted refinement. The critical success factors were: specifying physical constraints in the initial prompt, running validation on every iteration, and maintaining a tolerance buffer (500~m$^3$) for numerical feasibility. The project serves as a template for any constrained optimization problem in water resources engineering.

\end{document}